\chapter{Physics Without Objects}

\section{The Oldest Assumption in Physics}

Most physical and everyday reasoning shares an assumption old enough that it rarely gets stated outright: that reality is fundamentally made of things, substances, particles, objects, which persist through time and to which change happens. A wall is a thing; erosion happens to it. A person is a thing; aging happens to them. The thing comes first, ontologically, and change is something layered on top of it, an event a persisting substance undergoes rather than something constitutive of what the substance is. This book has never stated this assumption, because it never needed it, and it is worth pausing, now that the engineering is behind it, to notice what was quietly rejected along the way.

\section{What This Book Already Showed, Restated as Metaphysics}

Nothing here is new content. It is worth restating, plainly, what earlier chapters already established, because their consequence is more sweeping than their original engineering framing let on. Chapter 16 showed that an object's identity in this framework is not a label attached to a persisting substance, but a fixed point of the write cycle, a configuration the field's own dynamics keep reproducing. Chapter 20 refined this further: identity is not even a fixed point in the simplest sense, but a conserved trajectory through phase space, a closed orbit the system's dynamics maintain rather than a static thing sitting still. At no point in the intervening chapters was there ever occasion to introduce a substance, a persisting stuff that the wall, the fountain, or an institution was made of, over and above the trajectory itself. The wall was never a thing that happened to trace a stable trajectory. The wall was the trajectory, stable enough and long-lived enough to be worth a name.

\section{Process Ontology, Named Directly}

Chapter 6.6 used the phrase process ontology in passing, deliberately, without developing it, promising a return that this chapter now makes good on. What this book has built is a version of process ontology, the philosophical position, associated most closely with Alfred North Whitehead, that processes rather than substances are what is fundamental, and that things are best understood as patterns processes exhibit rather than as the bearers of those processes. It is worth being honest about the relationship between this book's own account and that tradition. This book did not set out to argue for process philosophy on its own philosophical terms, and it will not attempt, here, to adjudicate the centuries of argument that tradition has generated. What is genuinely interesting is the independence of the path that arrived at a structurally similar place: an engineering problem, keeping a generated wall consistent with itself across two observations, pushed, through nothing more exotic than the requirement that persistence actually be architectural rather than merely convenient, toward a formalism in which trajectories, not objects, were always doing the fundamental work. That a working piece of engineering reconstructs something with this shape, without having been built to prove any philosophical thesis, is itself a small piece of evidence worth taking seriously, whatever one ultimately makes of process philosophy as a general metaphysical position.

\section{What "Object" Still Means}

It would be a mistake to read this chapter as eliminativist, as though it argued that walls, fountains, and institutions do not really exist and are merely convenient fictions layered over an underlying flux. That is not this book's claim, and it is worth being precise about what the claim actually is. A dynamic fixed point, in the sense Chapters 16 and 20 developed, is a genuinely distinguished, non-arbitrary structure within the space of all possible trajectories: most trajectories are not fixed points, most configurations do not reproduce themselves under the write cycle's own dynamics, and the ones that do are picked out by real mathematical structure, not by an observer's convenience in naming them. The wall exists, in exactly the sense the word ordinarily carries. What has changed is only the analysis of what makes it exist as one continuing thing rather than a sequence of unrelated states: not a persisting substance underneath the trajectory, but the trajectory's own stability, the fact that it is a genuine fixed point rather than a passing configuration. Object remains a perfectly good word for what a stable, self-reproducing trajectory is. It is simply no longer doing the explanatory work substance metaphysics traditionally asked it to do.

\section{Why This Isn't Merely a Relabeling}

It is fair to ask whether anything beyond vocabulary actually turns on this distinction, and it is worth answering with a concrete case rather than an abstract assurance. Chapter 33.7 developed a genuine diagnostic from treating a robot's balance as a trajectory near a fixed point rather than as a binary state, upright or fallen: coherence, Φ, was shown to decline measurably before an observable loss of balance occurred, because the underlying claim being tracked was the stability of a trajectory, a continuous, quantifiable property, rather than the presence or absence of a persisting object called balance. A substance-first framing has no comparably natural place for this kind of leading indicator; balance, on that framing, is a state a robot either has or does not have, and there is no obvious quantity to track before the state changes. A trajectory-first framing predicts, structurally, that instability should be visible in advance, because instability is a fact about a trajectory's neighborhood, Chapter 19.5's local minima and saddle points, not merely about whether a threshold has yet been crossed. Chapter 36.5 found the same structure again at civilizational scale, an institution's coherence declining before its collapse becomes evident. This is not a difference in vocabulary. It is a difference in what kind of question the framework equips you to ask, and in at least these two cases, the trajectory-first question turned out to have a real, useful answer the substance-first question did not obviously have available.

\section{Looking Ahead}

If trajectories rather than objects are what is fundamental, a question immediately follows that this chapter has not yet addressed. Not every conceivable trajectory is one this book's architecture would permit; Chapter 18's constraint manifold C was built precisely to distinguish admissible configurations from merely imaginable ones. What, then, makes a trajectory real rather than merely coherent as an idea, and is admissibility itself doing more philosophical work here than this book has so far acknowledged? That is the question Chapter 38 takes up directly.
